\section*{Capítulo \ref{ch:b2:magma-emplacement} --- Os magmas e seu alojamento}

\begin{solution}{exo:b2:magma-emplacement:1}
$T = \qty{1073}{K}$. Para $w = 1$, $\ln w = 0$: $\log_{10}\eta = -3.545 +
9601/(1073 - 195.7) = -3.545 + 10.94 = 7.40$. Para $w = 4$, $\ln 4 = 1.386$:
$A = -3.545 + 1.155 = -2.390$, $B = 9601 - 3283 = 6318$, $C = 195.7 + 44.7 =
240.4$, e $\log_{10}\eta = -2.390 + 6318/832.6 = -2.390 + 7.588 = 5.20$. A
viscosidade cai por um fator $10^{7.40 - 5.20} = 10^{2.20} \approx 160$.
\end{solution}

\begin{solution}{exo:b2:magma-emplacement:2}
Riolito: $p_s = (6/0.44)^2 = \qty{186}{MPa}$, a uma profundidade de $\num{186e6}/(2600\times9.81)
= \qty{7.3}{km}$. Basalto: $p_s = (0.5/0.34)^2 = \qty{2.16}{MPa}$, a
$\num{2.16e6}/(2600\times9.81) = \qty{85}{m}$. O riolito começa a formar bolhas
no interior profundo da crosta, o basalto só perto da superfície (se não carregar CO$_2$).
\end{solution}

\begin{solution}{exo:b2:magma-emplacement:3}
$(1150 + 50)/2 = \qty{600}{\celsius}$ e $(1150 + 400)/2 = \qty{775}{\celsius}$. O
\omterm{def:b2:magma-emplacement:bodies}{dique} na rocha fria é resfriado mais bruscamente em seu contato, e tem a \omterm{def:b2:magma-emplacement:aureole}{margem resfriada}
mais larga.
\end{solution}

\begin{solution}{exo:b2:magma-emplacement:4}
(a) Uma \omterm{def:b2:magma-emplacement:bodies}{soleira}, paralela ao acamamento. (b) Um \omterm{def:b2:magma-emplacement:bodies}{dique}, que corta o acamamento.
(c) Um \omterm{def:b2:magma-emplacement:bodies}{batólito}, com mais de \qty{100}{km^2} e montado a partir de muitas intrusões.
(d) Um \omterm{def:b2:magma-emplacement:bodies}{lacólito}, que arqueou seu teto em domo. (e) Um \omterm{def:b2:magma-emplacement:bodies}{dique}: o granito não tem acamamento,
de modo que qualquer lâmina nele é um \omterm{def:b2:magma-emplacement:bodies}{dique}, qualquer que seja sua orientação.
\end{solution}

\begin{solution}{exo:b2:magma-emplacement:5}
$(1 - 1.35\phi)^{-2.5}$: para $\phi = 0.05$, $0.9325^{-2.5} = 1.19$, portanto
\qty{119}{Pa.s}, contra $1 + 2.5\times0.05 = 1.125$, \qty{113}{Pa.s}, com a fórmula de
Einstein (próximos, como deve ser no limite diluído). Para $0.20$: $0.73^{-2.5} =
2.20$, \qty{220}{Pa.s}; para $0.40$: $0.46^{-2.5} = 6.97$, \qty{697}{Pa.s}; para $0.60$:
$0.19^{-2.5} = 63.6$, cerca de \qty{6360}{Pa.s}. Travamento em $\phi = 1/1.35 = 0.74$.
\end{solution}

\begin{solution}{exo:b2:magma-emplacement:6}
$G = 300\times9.81 = \qty{2943}{Pa/m}$. (a) $\bar u = 2943\times1^2/(12\times100) =
\qty{2.45}{m/s}$; $q = \qty{2.45}{m^2/s}$; $Re = 2700\times2.45\times1/100 = 66$, laminar;
$t = 10000/2.45 = \qty{4080}{s}$, cerca de \qty{1.1}{h}. (b) $\bar u =
2943/(12\times\num{e6}) = \qty{2.45e-4}{m/s}$, $t = \qty{4.08e7}{s} = 472$ dias; com
$w = \qty{10}{m}$, $\bar u$ é cem vezes maior, \qty{0.0245}{m/s}, e $t =
\qty{4.08e5}{s} = 4.7$ dias. (c) O \omterm{def:b2:magma-emplacement:bodies}{dique} de \qty{1}{m} se resfria em cerca de
$0.5^2/\num{e-6} = \qty{2.5e5}{s}$, 2.9 dias; o \omterm{def:b2:magma-emplacement:bodies}{dique} de \qty{10}{m} em $5^2/\num{e-6} =
\qty{2.5e7}{s}$, 289 dias. O basalto (\qty{1.1}{h}) e o riolito no \omterm{def:b2:magma-emplacement:bodies}{dique}
largo (4.7 dias) chegam antes que seus \omterm{def:b2:magma-emplacement:bodies}{diques} se solidifiquem; o riolito no
\omterm{def:b2:magma-emplacement:bodies}{dique} de \qty{1}{m} (472 dias) não consegue.
\end{solution}

\begin{solution}{exo:b2:magma-emplacement:7}
Por quilograma de \omterm{def:b2:magma-emplacement:magma}{magma}, $n = 10/18.015 = \qty{0.555}{mol}$ de água e $nRT =
0.555\times8.314\times1173 = \qty{5414}{J}$; o fundido restante ocupa $0.99/2400 =
\qty{4.13e-4}{m^3}$. A \qty{100}{MPa}: $V = \qty{5.41e-5}{m^3}$, fração
$5.41/(5.41 + 41.3) = 0.12$. A \qty{10}{MPa}: $V = \qty{5.41e-4}{m^3}$, fração
$0.57$. A \qty{0.1}{MPa}: $V = \qty{0.0541}{m^3}$, fração $0.992$. O valor a
\qty{100}{MPa} é o menos confiável: o vapor de água a essa pressão é denso e
está longe de ser ideal (e boa parte da água ainda estaria, na verdade, dissolvida).
\end{solution}

\begin{solution}{exo:b2:magma-emplacement:8}
(a) $2300 + 60z = 2650$ dá $z = \qty{5.8}{km}$. (b) O fundido tem
\qty{2580}{kg/m^3}: $z = 280/60 = \qty{4.7}{km}$, cerca de \qty{1.2}{km} mais acima. (c)
Em $z = \qty{7.8}{km}$ a rocha tem $2300 + 60\times7.8 = \qty{2770}{kg/m^3}$; a
força é $(2770 - 2650)\times9.81 = \qty{1.2}{kN}$ por metro cúbico, para cima.
\end{solution}

\begin{solution}{exo:b2:magma-emplacement:9}
$v = 2\times300\times9.81\times(1000)^2/(9\times\num{e19}) = \qty{6.5e-11}{m/s}$, ou
$\num{6.5e-11}\times\num{3.156e7} = \qty{2.1e-3}{m}$ por ano, \qty{2.1}{mm/yr}. Subir
\qty{10}{km} leva $\num{e4}/\num{2.1e-3} = \qty{4.8}{Myr}$. A escala de tempo do resfriamento
é $(1000)^2/\num{e-6} = \qty{e12}{s} = \qty{32}{kyr}$: o corpo se solidifica depois de
subir cerca de $32000\times0.0021 \approx \qty{70}{m}$. Um \omterm{def:b2:magma-emplacement:ascent}{diápiro} magmático não consegue atravessar
uma crosta assim; ele precisa de encaixantes mais quentes e mais moles.
\end{solution}

\begin{solution}{exo:b2:magma-emplacement:10}
(a) $x_{\mathrm{OH}} = \sqrt{K_2x_m} = \sqrt{K_1K_2f}$, logo $x = K_1f + \frac12\sqrt{K_1K_2f}$.
A razão entre o primeiro termo e o segundo é $2\sqrt{K_1f/K_2}$: ela tende a zero
quando $f \to 0$ (a raiz quadrada domina) e cresce sem limite com $f$ (o
termo linear domina). (b) $K_1 f = \frac12\sqrt{K_1K_2 f}$ dá $K_1^2f^2 =
K_1K_2f/4$, logo $f = K_2/(4K_1)$. (c) $f = 0.2/(\num{4e-4}) = \qty{500}{MPa}$. A
\qty{100}{MPa}: $x_m = 0.0100$, $\frac12\sqrt{\num{e-4}\times0.2\times100} =
\frac12\times0.0447 = 0.0224$, logo $x = 0.0324$ e a fração em hidroxila é
$0.0224/0.0324 = \qty{69}{\percent}$.
\end{solution}

\begin{solution}{exo:b2:magma-emplacement:11}
(a) $1000/\num{4e6} = \qty{2.5e-4}{km^3}$ por ano, isto é, $\qty{2.5e5}{m^3}$ por ano,
ou $\num{2.5e5}/\num{3.156e7} = \qty{0.0079}{m^3/s}$, cerca de oito litros por segundo.
(b) O \omterm{def:b2:magma-emplacement:bodies}{dique} transporta $1000\times2\times0.8 = \qty{1600}{m^3/s}$, e fornece
$\qty{e12}{m^3}$ em \qty{6.25e8}{s}, cerca de 20 anos: $\num{5e-6}$ dos
\qty{4}{Myr}. (c) Cada lâmina contém $20\times\num{e4}\times\num{e4} = \qty{2e9}{m^3} =
\qty{2}{km^3}$, logo 500 lâminas, uma a cada $\num{4e6}/500 = 8000$ anos em média.
(d) $t = (10)^2/(4\times0.65^2\times\num{0.7e-6}) = \qty{8.5e7}{s}$, 2.7 anos. Cada
lâmina se solidifica milhares de anos antes que chegue a seguinte: a menos que o suprimento
venha em pulsos rápidos, o \omterm{def:b2:magma-emplacement:bodies}{plúton} nunca é um único grande corpo de \omterm{def:b2:magma-emplacement:magma}{magma}, como sugere a
datação dos grandes \omterm{def:b2:magma-emplacement:bodies}{plútons}.
\end{solution}

\begin{solution}{exo:b2:magma-emplacement:12}
(a) Com $x$ fixo, $\ln(T - T_0) = \text{const} - \frac12\ln t - x^2/(4\kappa t)$, cuja
derivada $-1/(2t) + x^2/(4\kappa t^2)$ se anula em $t^* = x^2/(2\kappa)$. Então
$e^{-x^2/(4\kappa t^*)} = e^{-1/2}$ e $\sqrt{4\pi\kappa t^*} = x\sqrt{2\pi}$, logo $\Delta
T_{\max} = (T_m - T_0)\,b\,e^{-1/2}/(x\sqrt{2\pi}) = (T_m - T_0)\,b/(x\sqrt{2\pi e})$. (b)
$\sqrt{2\pi e} = 4.13$: $x = 1100\times2/(100\times4.13) = \qty{5.3}{m}$, atingidos após
$t^* = 5.3^2/(2\times\num{0.7e-6}) = \qty{2.0e7}{s}$, cerca de 230 dias (a apenas 2.7
larguras de \omterm{def:b2:magma-emplacement:bodies}{dique}, a aproximação de fonte plana é grosseira). (c) O aquecimento decresce
como $b/x$: um \omterm{def:b2:magma-emplacement:bodies}{dique} delgado carrega pouco calor, e a rocha só é aquecida de modo apreciável
a poucas larguras de \omterm{def:b2:magma-emplacement:bodies}{dique}.
\end{solution}

\begin{solution}{pb:b2:magma-emplacement:1}
\textbf{1.} $\log_{10}\eta = -4.55 + 5510/(1423 - 611) = -4.55 + 6.786 = 2.24$;
$\eta = \qty{172}{Pa.s}$.

\textbf{2.} $-4.55 + 5510/762 = -4.55 + 7.231 = 2.68$, $\eta = \qty{480}{Pa.s}$: um
fator de $480/172 = 2.8$ para \qty{50}{K}.

\textbf{3.} $172\times(1 - 1.35\times0.20)^{-2.5} = 172\times0.73^{-2.5} = 172\times2.20 =
\qty{378}{Pa.s}$.

\textbf{4.} $\phi = 1/1.35 = 0.74$. As evidências de campo situam o fim do escoamento em massa
em uma fração de fundido de \qtyrange{30}{50}{\percent}, isto é,
\qtyrange{50}{70}{\percent} de cristais: antes, porque os cristais reais não são
esferas de um só tamanho.

\textbf{5.} Embasamento: $2800 - 2710 = \qty{90}{kg/m^3}$, o \omterm{def:b2:magma-emplacement:magma}{magma} é mais leve.
Sedimentos: $2450 - 2710 = \qty{-260}{kg/m^3}$, o \omterm{def:b2:magma-emplacement:magma}{magma} é mais denso.

\textbf{6.} $p = 2450\times9.81\times3000 = \qty{72.1}{MPa}$.

\textbf{7.} $90\times9.81 = \qty{883}{Pa/m}$.

\textbf{8.} $\bar u = 883\times2.0^2/(12\times378) = 3532/4536 = \qty{0.78}{m/s}$.

\textbf{9.} $q = 0.78\times2.0 = \qty{1.56}{m^2/s}$; ao longo de \qty{1.0}{km}, cerca de
\qty{1560}{m^3/s}.

\textbf{10.} $Re = 2710\times0.78\times2.0/378 = 11$: laminar.

\textbf{11.} $12000/0.78 = \qty{15400}{s}$, cerca de \qty{4.3}{h}.

\textbf{12.} $(1.0)^2/\num{0.7e-6} = \qty{1.43e6}{s}$, 16.5 dias, muito mais que
\qty{4.3}{h}: o \omterm{def:b2:magma-emplacement:bodies}{dique} permanece aberto enquanto o \omterm{def:b2:magma-emplacement:magma}{magma} escoa.

\textbf{13.} A \qty{3.0}{km} o \omterm{def:b2:magma-emplacement:magma}{magma} encontra sedimentos mais leves que ele: a
base dos sedimentos é seu \omterm{def:b2:magma-emplacement:ascent}{nível de flutuabilidade neutra}. Acima dela o \omterm{def:b2:magma-emplacement:magma}{magma} seria
mais denso que suas encaixantes, de modo que se espalha lateralmente ao longo do contato como uma
\omterm{def:b2:magma-emplacement:bodies}{soleira}.

\textbf{14.} $p_s = (0.40/0.34)^2 = \qty{1.38}{MPa}$, atingida a
$\num{1.38e6}/(2450\times9.81) = \qty{58}{m}$ abaixo da superfície.

\textbf{15.} $0.34\times\sqrt{72.1} = 0.34\times8.49 = \qty{2.9}{\percent}$, muito mais que
os \qty{0.40}{\percent} presentes: o \omterm{def:b2:magma-emplacement:magma}{magma} da \omterm{def:b2:magma-emplacement:bodies}{soleira} não exsolve água.

\textbf{16.} O CO$_2$ é muito menos solúvel que a água: ele satura e forma bolhas
logo no início da ascensão, em profundidade, de modo que o primeiro gás de um basalto é sobretudo CO$_2$.

\textbf{17.} Por quilograma, $n = 4.0/18.015 = \qty{0.222}{mol}$, $V = 0.222\times8.314\times1423/
\num{e5} = \qty{0.0263}{m^3}$; o fundido ocupa $0.996/2710 = \qty{3.68e-4}{m^3}$; a
fração é $0.0263/(0.0263 + 0.000368) = 0.986$, \qty{98.6}{\percent}.

\textbf{18.} $(1150 + 100)/2 = \qty{625}{\celsius}$.

\textbf{19.} Após 10 anos, $2\sqrt{\kappa t} = 2\sqrt{\num{0.7e-6}\times\num{3.156e8}} =
\qty{29.7}{m}$. O ponto está a $x = \qty{60}{m}$ do meio: $T = 100 +
525\,[\operatorname{erf}(-10/29.7) + \operatorname{erf}(110/29.7)] = 100 + 525\,(-0.365 +
1.000) = \qty{433}{\celsius}$.

\textbf{20.} No meio, $T = 100 + 1050\operatorname{erf}(50/2\sqrt{\kappa t})$: após
10 anos, $1050\operatorname{erf}(1.68) = 1050\times0.983$, $T = \qty{1132}{\celsius}$; após
50 anos, $2\sqrt{\kappa t} = \qty{66.5}{m}$, $1050\operatorname{erf}(0.752) = 1050\times0.713$,
$T = \qty{848}{\celsius}$.

\textbf{21.} Em $x = \qty{130}{m}$: \qty{294}{\celsius} após 300 anos,
\qty{295}{\celsius} após 363 anos (o máximo), \qty{295}{\celsius} após 420
anos. Sem calor latente, a rocha \qty{80}{m} acima da \omterm{def:b2:magma-emplacement:bodies}{soleira} nunca chega a atingir
\qty{300}{\celsius}.

\textbf{22.} $\num{9e8}\times1.772/(\num{2.5e6}\times1050) = 0.608$.

\textbf{23.} $\lambda = 0.64$: $e^{-0.4096}/(0.64\times1.6346) = 0.6639/1.0461 = 0.635$;
$0.65$: $0.6554/1.0673 = 0.614$; $0.66$: $0.6469/1.0886 = 0.594$. O valor $0.608$
fica entre $0.65$ e $0.66$, perto de $0.65$: $\lambda = 0.65$.

\textbf{24.} $100 + 1050/(1 + 0.642) = 100 + 640 = \qty{740}{\celsius}$. O calor
latente liberado na frente de solidificação mantém o lado do \omterm{def:b2:magma-emplacement:magma}{magma} a
\qty{1150}{\celsius} por mais tempo, de modo que o contato recebe mais calor.

\textbf{25.} $(50)^2/\num{0.7e-6} = \qty{3.57e9}{s}$, 113 anos.

\textbf{26.} $t = (100)^2/(16\times0.65^2\times\num{0.7e-6}) = \num{e4}/\num{4.73e-6} =
\qty{2.11e9}{s}$, logo $\num{2.11e9}/\num{3.156e7}$: a \omterm{def:b2:magma-emplacement:bodies}{soleira} está sólida após
\textbf{\qty[text-series-to-math=true]{67}{} anos}. É mais curto que
$(b/2)^2/\kappa$ pelo fator $1/(4\lambda^2) = 0.59$: essa escala de tempo só fixa
a ordem de grandeza, e as duas paredes da \omterm{def:b2:magma-emplacement:bodies}{soleira} extraem calor para uma encaixante fria ao
mesmo tempo, depressa o bastante para remover também o calor latente.
\end{solution}
